%%%%%%%%%%%%%%%%%%%%%%%%%%%%%%%%%
%%% CV Aparnna Hariharan 05/2024 %%%%%%
%%%%%%%%%%%%%%%%%%%%%%%%%%%%%%%%%

\PassOptionsToPackage{dvipsnames}{xcolor}
\documentclass[10pt,a4paper]{altacv}

% \newcommand{\mainspacing}{\vspace{8em}}

%Layout
% \geometry{left=1cm,right=9cm,marginparwidth=6.8cm,marginparsep=1.2cm,top=1.25cm,bottom=1.25cm,footskip=2\baselineskip}
\geometry{left=1cm,right=9cm,marginparwidth=7cm,marginparsep=1cm,top=1cm,bottom=1cm,footskip=2\baselineskip}

%Packages
\usepackage[utf8]{inputenc}
\usepackage[T1]{fontenc}
\usepackage[default]{lato}
\usepackage{hyperref} 
\usepackage{setspace}

%Colors

\definecolor{accent}{HTML}{000e17}
\definecolor{heading}{HTML}{000e17}
\definecolor{emphasis}{HTML}{696969}
\definecolor{body}{HTML}{01415f}

\colorlet{heading}{heading}
\colorlet{accent}{accent}
\colorlet{emphasis}{emphasis}
\colorlet{body}{body}

\renewcommand{\itemmarker}{{\small\textbullet}}
\renewcommand{\ratingmarker}{\faCircle}

%

\begin{document}
% \onehalfspacing
\setstretch{1.2}
\name{Aparnna Hariharan}
\tagline{}
\personalinfo{
    \homepage{\href{https://aparnnah.github.io}{aparnnah.github.io}}
    \mailaddress{aparnnah10@gmail.com}
    \github{\href{https://github.com/aparnnah}{github.com/aparnnah}}
    \linkedin{\href{https://linkedin.com/in/aparnna-hariharan}{linkedin.com/in/aparnna-hariharan}}
}

%

\begin{fullwidth}
\makecvheader
\end{fullwidth}

%

%%%%%%%%%%%%%%%%%%%%%%%%%%%%%%% Experience %%%%%%%%%%%%%%%%%%%%%%%%%%%%%%%

\cvsection[page1sidebar]{Experience}

\cvevent{Applied AI Evaluation Engineer}{Outlier}{August 2024 - August 2025}{}

\begin{itemize}
% \item Train generative artificial intelligence models, enhancing their coding generation capabilities.
% \item Evaluate the quality of AI-generated code, providing human-readable summaries and rationales for evaluations.
% \item Solve diverse coding challenges and write functional, efficient code while maintaining attention to detail.

% \item Trained AI models to enhance coding generation and accuracy
% \item Evaluated AI-generated code quality, delivering concise summaries and justifications to facilitate continuous improvements
% \item Developed efficient solutions to various coding challenges, rewriting code to enhance functionality and clarity

\item Trained and fine-tuned AI models to improve code generation quality, accuracy, and reliability.
\item Ran evaluations on AI-generated code across multiple LLMs, analyzing correctness, efficiency, and edge-case performance
\item  Wrote reference implementations from prompts and compared them against LLM-generated code
\item Identified compilation, logic, and performance errors and provided feedback to improve model output
\end{itemize}

%%%%%%%%%%%%%
%Experience 1
%%%%%%%%%%%%%

\cvevent{Software Developer Internship}{Ingram Micro}{January 2023 - August 2023}{}

\begin{itemize}
    \item Enhanced Azure monitor alerts by simplifying Azure log navigation and generating reports via terminal, using Python scripts with KQL queries, incorporating regex for simplified data parsing and additional user details for easier issue resolution
    % \item Developed Python scripts using KQL queries to enhance Azure monitor alerts, automate log navigation, and generate reports, with regex for simplified data parsing and additional user details for easier issue resolution
    \item Developed Python scripts to produce customizable charts, providing quick access to valuable insights and trend analysis for multiple team asset maintenance, using Jira time log pull functionality
    \item Designed and developed a product analytic report retrieval system integrated into an existing analytics API
    % \item Implemented tool-time production usage tracking for Freshdesk, enabling the collection of previously unavailable data on feature usage over the past 6 months using KQL 
    % \item Implemented a Base64 encoding and decoding feature and a Java-based random generator to the tool-time website

\end{itemize}

\medskip

%%%%%%%%%%%%%
%Experience 2
%%%%%%%%%%%%%

\cvevent{Computer Science Research Assistant}{Ontario Tech University}{May 2020 - August 2020}{}

\begin{itemize}
    \item Designed and implemented a traffic manager in C++ that can change and control the vehicles in Carla (an open source connected vehicle simulation) 
    \item Developed a scenario runner in Python that allows different traffic scenarios to be executed in the Carla simulation
    \item Wrote an analysis report on connected vehicles and the design of the traffic manager implementation on Carla

\end{itemize}

\medskip

%%%%%%%%%%%%%%%%%%%%%%%%%%%%%%% Projects %%%%%%%%%%%%%%%%%%%%%%%%%%%%%%%

\cvsection[page2sidebar]{Projects}


%%%%%%%%%%%%%
%Project 1
%%%%%%%%%%%%%

% \medskip

% \cvproject{\href{https://aparnnah.github.io/socialmediausage.html}{Social Media Usage Report | R}}{}{}

% \begin{itemize}
%   \item Analyzed publicly available data on global screen time and social media habits, highlighting variations in age groups and platform preferences through visuals
% \end{itemize}

% \begin{itemize}
%   \item Explored global screen time and social media habits, highlighting variations in age groups and platform preferences through visuals
% \end{itemize}

%%%%%%%%%%%%%
%Project 4
%%%%%%%%%%%%%


% \cvproject{Find Carmen Sandiego}{C\#}
% {\href{https://www.github.com/aparnnaH/FindCarmenSandiego}{github.com/aparnnaH/FindCarmenSandiego}}

% \begin{itemize}
%   \item {Designed and implemented a windows game where the players need to solve a series of 8 arcade puzzles in order to stop Carmen Sandiego}
% \end{itemize}
% %%%%%%%%%%%%
% Project 3
% %%%%%%%%%%%%
% \medskip

% \cvproject{\href{https://github.com/aparnnaH/battleSystem}{{Battle Simulator | Java}}}

% \begin{itemize}
%   \item Developed a battle system that will simulate a battle between a user and a computer AI 
% \end{itemize}

%%%%%%%%%%%%%
%Project 5
%%%%%%%%%%%%%

% \medskip

% \cvproject{\href{https://github.com/taylor-young0/Object-Removal}{Object Removal | Python}}{Python}{}

% \cvproject{\href{https://github.com/taylor-young0/Object-Removal}{Object Removal | Python}}{}
% \begin{itemize}
%   \item {Developed an interactive image editing tool to remove object and generate a fill based on surrounding pixels.}
% \end{itemize}

%%%%%%%%%%%%%
%Project 1
%%%%%%%%%%%%%

% \medskip

\cvproject{\href{https://github.com/aparnnaH/Twitter-Emoji-Prediction-Project}{Twitter Emoji Prediction | Python}}{Repository}{}

\begin{itemize}
  \item Developed a deep learning NLP pipeline for multi-class emoji prediction using LSTM, GRU, and TextCNN with text preprocessing and sequence modeling.
\end{itemize}

%%%%%%%%%%%%%
%Project 3
%%%%%%%%%%%%%
% \medskip

% \cvproject{\href{https://github.com/aparnnaH/battleSystem}{{Battle Simulator | Java}}}

% \begin{itemize}
%   \item Developed a battle system that will simulate a battle between a user and a computer AI 
% \end{itemize}

%%%%%%%%%%%%%
%Project 1
%%%%%%%%%%%%%

% \medskip

% \cvproject{\href{https://github.com/aparnnaH/recipe}{Recipe Website | Django}}{Repository}{}

% \begin{itemize}
%   \item Designed and implemented a recipe web app on Django that allows users to browse, post, comment, categorize and favourite recipes
% \end{itemize}

% \medskip

\cvproject{\href{https://github.com/aparnnaH/Heart-Disease-Unsupervised-Learning/tree/main}{Heart Disease Unsupervised Learning | Python}}{Repository}{}

\begin{itemize}
  \item Developed a pipeline using dimensionality reduction and visualization to uncover risk-based subgroups and clinical patterns in heart disease.
\end{itemize}


%%%%%%%%%%%%%
%Project 2
%%%%%%%%%%%%%
% \medskip

% \cvproject{\href{https://www.github.com/aparnnaH/todoListApp}{To-do List App | Swift}}{Repository}{}
% \begin{itemize}
%   \item Designed and implemented an IOS app that allows the user to keep track of different tasks with various deadlines
% \end{itemize}

\end{document}